\documentclass[11pt,a4paper]{article}

\usepackage[utf8]{inputenc}
\usepackage[T1]{fontenc}
\usepackage{lmodern}
\usepackage[provide=*,spanish]{babel}
\usepackage[a4paper,margin=2.35cm,headheight=15pt]{geometry}
\usepackage{microtype}
\usepackage{booktabs}
\usepackage{tabularx}
\usepackage{longtable}
\usepackage{array}
\usepackage{xcolor}
\usepackage{hyperref}
\usepackage{xurl}
\usepackage{enumitem}
\usepackage{listings}
\usepackage{fancyhdr}
\usepackage{parskip}
\usepackage{needspace}
\usepackage{graphicx}
\usepackage{pdfpages}
\usepackage{pgf-umlsd}

\definecolor{ucaldasink}{HTML}{2D3142}
\definecolor{ucaldasslate}{HTML}{4F5D75}
\definecolor{ucaldasaccent}{HTML}{BF4520}
\definecolor{ucaldasrule}{HTML}{D8DADF}

\hypersetup{
  colorlinks=true,
  linkcolor=ucaldasink,
  urlcolor=ucaldasaccent,
  citecolor=ucaldasink,
  pdftitle={Gestión de Monitores Académicos},
  pdfauthor={Equipo del proyecto}
}

\setlength{\parindent}{0pt}
\setlength{\parskip}{0.55em}
\setlist[itemize]{leftmargin=1.5em,itemsep=0.2em,topsep=0.25em}
\setlist[enumerate]{leftmargin=1.7em,itemsep=0.25em,topsep=0.25em}
\renewcommand{\arraystretch}{1.2}
\newcolumntype{Y}{>{\raggedright\arraybackslash}X}

\lstdefinestyle{comandos}{
  basicstyle=\ttfamily\small,
  breaklines=true,
  breakatwhitespace=false,
  columns=fullflexible,
  keepspaces=true,
  showstringspaces=false,
  frame=single,
  rulecolor=\color{ucaldasrule},
  backgroundcolor=\color{black!3},
  xleftmargin=0.4em,
  xrightmargin=0.4em,
  aboveskip=0.6em,
  belowskip=0.6em,
  literate=
    {á}{{\'a}}1 {é}{{\'e}}1 {í}{{\'i}}1 {ó}{{\'o}}1 {ú}{{\'u}}1
    {Á}{{\'A}}1 {É}{{\'E}}1 {Í}{{\'I}}1 {Ó}{{\'O}}1 {Ú}{{\'U}}1
    {ñ}{{\~n}}1 {Ñ}{{\~N}}1 {ü}{{\"u}}1 {Ü}{{\"U}}1
}
\lstset{style=comandos}

\pagestyle{fancy}
\fancyhf{}
\fancyhead[L]{\small Gestión de Monitores Académicos}
\fancyhead[R]{\small Universidad de Caldas}
\fancyfoot[L]{\small Ingeniería de Software III · 2026-2}
\fancyfoot[R]{\small Página \thepage}
\renewcommand{\headrulewidth}{0.4pt}
\renewcommand{\footrulewidth}{0.4pt}

\newcommand{\portada}[2]{%
  \hypersetup{pageanchor=false}
  \begin{titlepage}
    \centering
    \vspace*{2.5cm}
    {\large Universidad de Caldas\par}
    \vspace{0.5cm}
    {\large Ingeniería de Software III · 2026-2\par}
    \vspace{2.5cm}
    {\Huge\bfseries Gestión de Monitores Académicos\par}
    \vspace{0.7cm}
    {\LARGE #1\par}
    \vspace{0.7cm}
    {\large #2\par}
    \vfill
    {\large Versión 1.0 · 20 de noviembre de 2026\par}
    \vspace{0.4cm}
    {\normalsize Equipo del proyecto\par}
  \end{titlepage}
  \hypersetup{pageanchor=true}
  \pagenumbering{roman}
}

\newcommand{\codigo}[1]{\texttt{#1}}
\newcommand{\archivo}[1]{\path{#1}}
\newcommand{\componente}[1]{\texttt{#1}}


\begin{document}

\portada{Manual de usuario}{Guía de acceso y gestión de monitores}
\tableofcontents
\clearpage
\pagenumbering{arabic}

\section{Antes de comenzar}

El sistema permite registrar, consultar, editar y eliminar monitores académicos. Para usarlo, se necesita que el equipo donde corre la aplicación tenga Docker Desktop y Docker Compose v2 activos.

\subsection{Iniciar la aplicación}

Desde la carpeta del proyecto, abrir una terminal y ejecutar:

\begin{lstlisting}
docker compose up --build -d
docker compose ps
\end{lstlisting}

Esperar a que los cinco servicios indiquen estado saludable y abrir \url{http://localhost:8080} en el navegador. La primera ejecución puede tardar varios minutos.

Para cerrar la aplicación sin borrar la información, ejecutar:

\begin{lstlisting}
docker compose down
\end{lstlisting}

\section{Reconocer la pantalla}

La pantalla principal contiene:
\begin{itemize}
  \item un indicador de conexión con el backend;
  \item el formulario para registrar o editar;
  \item el buscador por código y el botón \textbf{Ver todos};
  \item la tabla de monitores con las acciones \textbf{Editar} y \textbf{Eliminar};
  \item mensajes de resultado o de error.
\end{itemize}

Las materias se cargan desde el sistema. No se crean ni modifican desde esta pantalla.

\section{Registrar un monitor}

\begin{enumerate}
  \item En el formulario \textbf{Registrar monitor}, escribir un código único. Se aceptan letras, números y guiones, hasta 20 caracteres; el sistema lo guarda en mayúsculas. Ejemplo: \codigo{M-008}.
  \item Escribir el nombre completo, con entre 3 y 150 caracteres.
  \item Ingresar un correo válido, de hasta 150 caracteres.
  \item Seleccionar una materia de la lista.
  \item Presionar \textbf{Guardar}.
\end{enumerate}

Si el registro se completa, aparece un mensaje de confirmación y la tabla se actualiza. Los campos obligatorios deben estar completos. Un código ya utilizado no se puede reutilizar.

\section{Buscar y consultar}

\subsection{Ver todos los registros}

Presionar \textbf{Ver todos} para limpiar la búsqueda y volver a mostrar la lista completa. Si todavía no hay registros, la tabla indica que está vacía.

\subsection{Buscar por código}

Escribir el código en \textbf{Buscar por código} y presionar \textbf{Buscar}. La búsqueda no distingue mayúsculas de minúsculas; por ejemplo, \codigo{m-001} encuentra \codigo{M-001}. Para regresar a la lista completa, presionar \textbf{Ver todos}.

\section{Editar un monitor}

\begin{enumerate}
  \item Encontrar el registro en la tabla.
  \item Presionar \textbf{Editar}. Sus datos se cargan en el formulario.
  \item Modificar nombre, correo o materia.
  \item Presionar \textbf{Guardar cambios}.
\end{enumerate}

El código identifica al monitor y no se puede cambiar. Para descartar los cambios, presionar \textbf{Cancelar edición}; el formulario vuelve al modo de registro sin guardar.

\section{Eliminar un monitor}

\begin{enumerate}
  \item Presionar \textbf{Eliminar} en la fila del monitor.
  \item Revisar el código y confirmar en el diálogo.
  \item Presionar \textbf{Eliminar} para continuar o \textbf{Cancelar} para conservar el registro.
\end{enumerate}

La eliminación es permanente dentro de la base de datos activa y no se puede deshacer desde la interfaz. Cerrar el diálogo, pulsar \codigo{Esc} o hacer clic fuera de él cancela la operación.

\section{Mensajes y qué hacer}

\begin{tabularx}{\textwidth}{p{0.38\textwidth}Y}
\toprule
\textbf{Mensaje o situación} & \textbf{Qué revisar} \\
\midrule
Código obligatorio o formato inválido & Escribir de 1 a 20 letras, números o guiones. \\
Nombre obligatorio o longitud inválida & Ingresar entre 3 y 150 caracteres. \\
Correo inválido & Comprobar que tenga usuario, \codigo{@} y un dominio con punto; no usar espacios. \\
Materia obligatoria & Seleccionar una opción de la lista. \\
Ya existe un monitor con ese código & Elegir un código diferente. \\
No existe un monitor con ese código & Volver a la lista y confirmar que el registro todavía exista. \\
No se pudo conectar con el servidor & Revisar que Docker esté en ejecución y que el frontend y los servicios aparezcan saludables. \\
No se pudieron cargar materias o monitores & Revisar la conexión y solicitar ayuda técnica si el problema persiste. \\
\bottomrule
\end{tabularx}

Los errores de validación del sistema se muestran tal como los genera el backend. Si una operación no puede completarse, corregir los datos o contactar a quien administra la instalación.

\section{Datos y privacidad}

El formulario almacena código, nombre, correo y materia. El sistema no solicita inicio de sesión: cualquier persona con acceso a la dirección de la aplicación puede operar sobre los registros. No publicar el servicio en una red abierta sin controles adicionales.

Los datos permanecen al cerrar y volver a iniciar los contenedores. Solo se eliminan al borrar explícitamente el volumen de base de datos o reinicializarla.

\section{Ayuda técnica}

\subsection{La aplicación no abre}

Comprobar que Docker Desktop esté abierto y que los puertos necesarios estén libres: 8080, 8081, 8082, 8083 y 5432. Si Docker informa que un puerto está ocupado, contactar a quien administra el entorno para cambiar el puerto en la configuración.

\subsection{Un servicio no está saludable}

Una persona con acceso a la carpeta del proyecto puede revisar:

\begin{lstlisting}
docker compose ps
docker compose logs -f <servicio>
\end{lstlisting}

Los servicios son \codigo{database}, \codigo{python-service}, \codigo{php-service}, \codigo{java-controller} y \codigo{frontend}. El orden de inicio importa: un servicio puede quedar esperando mientras la capa anterior no esté lista.

\Needspace{8\baselineskip}
\subsection{Advertencia de pérdida de datos}

\textbf{No ejecutar \codigo{docker compose down -v} ni \codigo{./scripts/reset-db.sh} si se quieren conservar los registros.}

Estos comandos eliminan el volumen de PostgreSQL y vuelven a cargar los datos de demostración.

\section{Documentos relacionados}

\begin{itemize}
  \item \archivo{manual-tecnico.pdf}: despliegue, arquitectura y diagnóstico.
  \item \archivo{manual-tecnico.tex}: fuente LaTeX del manual técnico.
  \item \archivo{arquitectura.pdf}: diagrama de la aplicación.
  \item \archivo{api-contract.md}: contrato de rutas y respuestas.
\end{itemize}

\end{document}
